\section{Conclusiones}

\subsection{Síntesis del sistema construido}

El trabajo entrega un sistema de gestión de acceso vehicular ensamblado a partir
de dos subsistemas con responsabilidades separadas. El subsistema de visión
encadena detección de vehículos sobre pesos preentrenados, seguimiento
multiobjeto, un detector de placas entrenado para este dominio, un reconocedor de
matrículas con arquitectura ResNet18--BiLSTM--CTC y dos clasificadores de
atributos. El subsistema de gestión expone diecisiete rutas HTTP ---operativas, de consulta,
de autenticación y de diagnóstico--- y un canal persistente de notificación sobre
FastAPI, PostgreSQL y una cola de tareas, y ofrece un cliente web con seis vistas. El despliegue se resuelve con contenedores,
exportación a formatos intermedios y compilación de motores especializados para el
acelerador de destino.

La magnitud del artefacto es verificable y no estimada: el sistema comprende $62$
puntos de función y $6.946$ líneas de producto distribuidas en $88$ archivos, a
las que se suman $3.537$ líneas de pruebas y $6.225$ de herramientas de datos. El
conjunto de datos que lo sustenta también es verificable: $496$ matrículas reales
confirmadas por un operador, $359.992$ matrículas sintéticas con formato
boliviano, $4.006$ recortes de vehículo y $868$ recortes de placa con caja
revisada manualmente.

Esa verificabilidad alcanza al software. El subsistema de gestión se encuentra
integrado con el motor de visión: la tarea de inferencia decodifica el fotograma,
invoca al motor, agrega las detecciones por trayectoria y persiste el vehículo, la
sesión y la observación de forma atómica, con la latencia de inferencia registrada
por fotograma. El esquema se materializa en tres revisiones de migración que se
aplican en orden, el cierre de sesión opera por emparejamiento exacto de matrícula
con umbral de confianza y admite el cierre manual auditado, y la aplicación web
ofrece seis vistas. La verificación automática comprende $140$ pruebas unitarias y
$4$ de integración que se ejecutan contra una base PostgreSQL real con las
migraciones aplicadas, con una cobertura del $88\%$ sobre el código de la
aplicación.

El límite que corresponde declarar es de otra naturaleza. El contraste de la
matrícula contra un padrón de vehículos habilitados, previsto como requerimiento
RF-10, no forma parte de la entrega: la entidad de vehículos registra los
vehículos observados en los accesos, no un padrón previo. Lo que tampoco se
realizó es la \textbf{validación de campo de extremo a extremo}: la cadena
completa con las cámaras del recinto y la simulación con las dos fuentes
sincronizadas. Esa validación es la que sostendría la medición de la determinación
de eventos ($H_{1.5}$) y la que permitiría afirmar que el sistema opera como un
todo continuo, y no solo que cada componente funciona y está probado. En
consecuencia, los resultados que se reportan corresponden a los componentes
efectivamente ejecutados y se atribuyen a ellos, no a una operación en condiciones
reales de recinto.

\subsection{Contrastación de la hipótesis general}

La hipótesis general exige que el sistema alcance simultáneamente los criterios
de desempeño establecidos para la detección e identificación de vehículos, el
reconocimiento de placas y la determinación de eventos. Por su forma conjuntiva,
basta que un solo término no se cumpla para que la hipótesis no se respalde, y
basta que un término no se haya medido para que no pueda declararse verificada.

El estado de cada término se resume en la Tabla~\ref{tab:contraste_hipotesis}:
dos criterios se cumplen (identificación de tipo y detección de placas), tres no
se alcanzan (color predominante, exactitud por carácter y coincidencia exacta) y
dos no fueron medidos (detección de vehículos y determinación de eventos).

\textbf{La hipótesis general no se respalda.} Dos de sus componentes medidos
quedan por debajo del umbral comprometido, de modo que el sistema, en su estado
actual, no satisface el conjunto de criterios que la hipótesis exige. Se acepta
la hipótesis nula para esos componentes.

La conclusión debe formularse con una precisión adicional: la hipótesis nula
tampoco queda demostrada en su totalidad, porque dos de sus términos no fueron
sometidos a medición. Respecto de la detección de vehículos y de la determinación
de eventos, el trabajo no afirma que el sistema falle ni que acierte: afirma que
no lo sabe, y explicita por qué. Un pronunciamiento en cualquiera de las dos
direcciones carecería de respaldo experimental.

\subsection{Cumplimiento de los objetivos específicos}

La Tabla~\ref{tab:cumplimiento_objetivos} contrasta cada objetivo específico con
la evidencia producida y con el grado de cumplimiento que esa evidencia sostiene.
Se emplean tres estados: \emph{cumplido} cuando el objetivo se alcanzó y su
resultado es verificable; \emph{parcial} cuando el artefacto existe y fue
evaluado, pero algún criterio asociado no se alcanzó o alguna medición quedó
pendiente; y \emph{no cumplido} cuando no existe evidencia que lo respalde. En
este trabajo no se registra ningún objetivo en el tercer estado.

\begin{longtable}{
        P{0.30}
        P{0.12}
        P{0.50}
    }
    \caption{Cumplimiento de los objetivos específicos}
    \label{tab:cumplimiento_objetivos}\\

    \toprule
    \rowcolor{gray!20}
    \textbf{Objetivo específico} &
    \textbf{Estado} &
    \textbf{Evidencia y observación} \\
    \midrule
    \endfirsthead

    \multicolumn{3}{c}{\textit{Continuación de la Tabla
    \ref{tab:cumplimiento_objetivos}}} \\
    \toprule
    \rowcolor{gray!20}
    \textbf{Objetivo específico} &
    \textbf{Estado} &
    \textbf{Evidencia y observación} \\
    \midrule
    \endhead

    Analizar las condiciones de operación, los puntos de acceso y los
    requerimientos técnicos y funcionales. &
    Cumplido &
    Requerimientos funcionales, no funcionales y de infraestructura
    especificados con identificadores trazables; condiciones de captura,
    iluminación y ángulo caracterizadas a partir del material recolectado. \\
    \midrule

    Construir y preparar un conjunto de datos representativo de las condiciones
    reales. &
    Cumplido &
    $496$ matrículas reales verificadas por operador, $359.992$ sintéticas,
    $4.006$ recortes de vehículo y $868$ recortes de placa con caja revisada.
    Partición por sesión y por bloque horario, sin sesiones compartidas entre
    entrenamiento y prueba. El volumen real es la restricción dominante del
    módulo de reconocimiento y así se declara. \\
    \midrule

    Adaptar y entrenar modelos de detección e identificación de tipo y color, y
    de localización de placas. &
    Parcial &
    Tipo de vehículo: F1 macro $0{,}8887$ sobre el umbral de $0{,}85$. Color:
    F1 macro $0{,}8124$ bajo el umbral. Localización de placas: exhaustividad
    $0{,}9714$ y mAP@50 $0{,}9898$. Detección de vehículos: el modelo opera con
    pesos preentrenados sin ajuste fino y su mAP no fue medido por ausencia de
    cajas de referencia humana. \\
    \midrule

    Desarrollar y entrenar un modelo de reconocimiento óptico de caracteres para
    placas vehiculares y evaluarlo bajo distintas condiciones de visibilidad. &
    Parcial &
    Modelo construido y entrenado en dos etapas, evaluado sobre $101$ placas
    aisladas y sobre $88$ recortes de vehículo con texto verificado. Exactitud
    por carácter $0{,}789$ y $0{,}737$; coincidencia exacta $0{,}356$ y
    $0{,}261$. El criterio de $0{,}90$ de exactitud por carácter no se alcanza;
    la causa se analiza en la Subsección~\ref{subsec:brecha_ocr}. \\
    \midrule

    Integrar los modelos en un sistema de gestión con registro y consulta de
    eventos. &
    Cumplido &
    Modelo de datos con restricciones e índices, tres revisiones de migración,
    diecisiete rutas HTTP, canal persistente de notificación, cliente web con seis
    vistas y emparejamiento de salida con abstención.
    Verificado por $140$ pruebas unitarias y $4$ de integración contra una base
    real. La validación de extremo a extremo con las cámaras del recinto
    corresponde al objetivo de evaluación, no a este. \\
    \midrule

    Evaluar el desempeño del sistema y el efecto de las técnicas de
    optimización. &
    Parcial &
    Cinco de los siete términos de hipótesis fueron medidos sobre conjuntos de
    prueba independientes. La optimización fue aplicada y su efecto sobre la
    huella en disco está cuantificado (reducción del $75\%$ al pasar de precisión
    simple a media precisión), pero la latencia, la tasa de fotogramas y la
    degradación de exactitud atribuible a la cuantización no llegaron a medirse. \\
    \bottomrule
    \multicolumn{3}{r}{\footnotesize\textit{Fuente: Elaboración propia}} \\*
\end{longtable}

Los tres objetivos parciales comparten un mismo patrón: el artefacto se
construyó, se ejecutó y se midió, pero la medición no alcanzó el umbral
comprometido o no cubrió todos los términos previstos. Ninguno de los tres
constituye un objetivo abandonado. Se trata de objetivos cuyo cumplimiento está
acotado por condiciones que este apartado identifica y que las recomendaciones
traducen en trabajo concreto.

\subsection{La brecha del reconocimiento de matrículas}
\label{subsec:brecha_ocr}

Es el componente de mayor distancia entre lo comprometido y lo obtenido, y por
eso recibe un tratamiento propio. El criterio exigía, de forma simultánea,
exhaustividad de detección igual o superior al $90\%$ y exactitud por carácter
igual o superior al $90\%$, con una coincidencia exacta adicional del $80\%$. El
componente de detección se cumplió con holgura ($0{,}9714$ de exhaustividad y
$0{,}9898$ de mAP@50). El déficit es íntegramente de reconocimiento: $0{,}789$ de
exactitud por carácter sobre recortes aislados y $0{,}737$ sobre el sistema
ensamblado, con coincidencia exacta de $0{,}356$ y $0{,}261$ respectivamente.

La conclusión que sostiene la evidencia es que \textbf{la restricción dominante
no es la arquitectura del modelo, sino la calidad y el volumen del material
disponible}. Cuatro mediciones independientes convergen en esa dirección.

\paragraph{Primera causa: la resolución física del recorte.}
Un recorte de $40 \times 13$ píxeles que contiene siete caracteres dispone de
aproximadamente cinco píxeles de ancho por glifo. A esa escala la información no
está presente en la imagen: ningún modelo puede reconstruir un trazo que no fue
capturado. La degradación es monótona con la altura del recorte ---$0{,}605$ por
debajo de $20$ píxeles, $0{,}824$ entre $20$ y $30$, $0{,}803$ entre $30$ y $40$
y $0{,}847$ por encima de $40$---, lo que descarta una explicación por
variabilidad muestral y confirma una relación causal con la resolución.

Este dato tiene una consecuencia decisiva para interpretar la brecha: aun
eliminando por completo los recortes de baja resolución y evaluando solo el mejor
tramo, el sistema alcanza $0{,}847$, todavía por debajo del $0{,}90$ exigido. La
infraestructura de captura disponible no produce, en la generalidad de los casos,
la calidad de imagen que el criterio supone.

\paragraph{Segunda causa: el volumen de datos reales.}
El ajuste fino dispone de $395$ placas reales para un alfabeto de $37$ símbolos.
La curva de exactitud seguía aumentando al finalizar las $500$ épocas, lo que
indica que el modelo no había agotado la información disponible: el límite
aparece antes por cantidad de datos que por convergencia del entrenamiento. El
preentrenamiento sobre $359.992$ matrículas sintéticas elevó el piso del
desempeño, pero no su techo, porque persiste una diferencia de dominio medible:
la luminancia mediana de los recortes reales es de $76{,}6$ frente a $106{,}4$ en
los sintéticos, con desviación estándar de $20{,}7$ frente a $34{,}4$.

\paragraph{Tercera causa: el costo de la localización.}
La comparación entre evaluar con la caja provista por un operador y con la caja
producida por el detector propio aísla el error atribuible al encuadre:
$0{,}799$ frente a $0{,}737$, es decir $6{,}2$ puntos de exactitud por carácter
que no corresponden al reconocedor. El dato se confirma desde el otro extremo:
de las $700$ propuestas del etiquetador automático que sí correspondían a una
placa, solo $183$ ($26{,}1\%$) resultaron utilizables sin corrección. La
geometría del recorte, y no solo su contenido, condiciona la lectura.

\paragraph{Cuarta causa: el formato impide corregir por gramática.}
Se implementó y se midió un post-procesamiento determinista basado en la
gramática oficial de las matrículas bolivianas. Su efecto fue contrario al
esperado: redujo la exactitud por carácter de $0{,}789$ a $0{,}786$, porque el
formato admite letras dentro del bloque numérico y la regla convierte lecturas
correctas en incorrectas. El mecanismo se conserva en el código como referencia
sintáctica y queda desactivado por omisión. La lección es que el formato
boliviano ofrece menos restricciones explotables que otros formatos, y que la
corrección sintáctica no sustituye la calidad de captura.

\paragraph{Por qué no se trata de un problema de ajuste.}
Tres observaciones descartan que la brecha se cierre con calibración de
hiperparámetros. Primero, dos arquitecturas de concepción distinta ---el CRNN
entrenado para este dominio y un motor genérico de terceros--- convergen en la
misma franja de desempeño ($0{,}789$ y $0{,}601$), lo que apunta a un límite del
material y no del método. Segundo, la diferencia entre la mejor época observada
($0{,}793$) y la reportada ($0{,}789$) es inferior al error de muestreo del
conjunto de prueba, estimado en $\pm 5$ puntos con $n = 101$; elegir el máximo
habría introducido un sesgo de selección sin mejora real. Tercero, las
intervenciones que sí produjeron cambios de magnitud fueron todas de
preprocesamiento y de datos, no de arquitectura: la normalización de contraste
llevó la exactitud de $0{,}14$ a $0{,}676$, el aumento geométrico de recorte de
$0{,}594$ a $0{,}737$, y la corrección del criterio de selección de época de
$0{,}7937$ a $0{,}8124$ en el clasificador de color, sin modificar el modelo ni
una sola imagen.

\paragraph{Consecuencia operativa.}
El reconocedor emite siempre una cadena plausible y acumula aciertos parciales:
por eso supera al motor genérico en exactitud por carácter y queda por debajo en
coincidencia exacta, dado que el motor genérico devuelve vacío cuando no alcanza
su umbral de confianza. Para el control de acceso, esa diferencia no es neutra:
una matrícula incorrecta registrada es peor que una matrícula no leída, porque
traslada el error a la base de datos y a la trazabilidad del vehículo. De aquí se
desprende una conclusión de diseño, y el sistema la implementa: el emparejamiento
entre ingreso y salida exige coincidencia exacta de matrícula y un umbral de
confianza de $0{,}6$, y cuando no lo alcanza deja la sesión abierta en lugar de
cerrarla con la mejor aproximación. La alternativa ---usar el extractor del
clasificador de tipo como señal de identidad--- se midió y se descartó: sobre
$400$ recortes, el $19{,}9\%$ de los pares de vehículos distintos superó el umbral
de similitud de $0{,}95$, de modo que ningún punto de corte los separa. La
consecuencia operativa es que el sistema no opera de forma autónoma frente a una
lectura dudosa: la deja visible y la resuelve un operador con registro de motivo y
responsable. Esa es la respuesta correcta a un techo de reconocimiento medido, y
no un rodeo. 

\paragraph{Qué haría falta para cerrar la brecha.}
Las mediciones permiten formular la condición necesaria con una cifra: una altura
de placa igual o superior a $40$ píxeles en el punto de lectura, combinada con un
volumen de datos reales un orden de magnitud superior al actual. En esas
condiciones el sistema ya alcanza $0{,}847$ de exactitud por carácter sobre el
tramo de mayor resolución, y la brecha restante es abordable con las técnicas que
se recomiendan en el apartado siguiente ---votación sobre múltiples fotogramas de
la trayectoria, hoy limitada al mejor fotograma, y decodificación restringida
solo en los bloques donde la gramática no es ambigua---. Ninguna de esas técnicas
fue ejecutada en este trabajo: se enuncian como trabajo futuro y no como
resultado. El aporte de esta sección es haber establecido, con evidencia, cuál es
la condición de captura por debajo de la cual el criterio del $90\%$ es
inalcanzable con cualquier modelo.

\subsection{Aportes verificables}

\paragraph{Aporte teórico.}
La sistematización de fundamentos se complementa con un resultado empírico que no
estaba en el perfil: la demostración de que, en reconocimiento de matrículas, el
techo de desempeño está gobernado por la resolución de captura y el volumen de
datos antes que por la capacidad del modelo, y la estimación de la condición de
captura mínima compatible con un umbral del $90\%$.

\paragraph{Aporte práctico.}
El pipeline integrado, el conjunto de datos construido y verificado, el sistema
de gestión con su interfaz de programación y su cliente web, y los protocolos de
evaluación quedan disponibles como artefactos reutilizables. El conjunto de datos
es el activo de mayor valor: $496$ matrículas reales con texto confirmado por
operador y $868$ recortes de placa con caja revisada no se obtienen sin el
trabajo de inspección que aquí se documenta.

\paragraph{Aporte académico.}
El trabajo deja registro de resultados negativos que rara vez se publican: el
post-procesamiento sintáctico medido y desactivado, las configuraciones de
entrenamiento descartadas con su efecto cuantificado, y las dos hipótesis que no
pudieron medirse con su causa. Ese registro evita que un trabajo posterior repita
el mismo camino y permite discutir el desempeño con datos en lugar de
impresiones.
